% =====================================================================
\section{Real deployments}
\label{sec:casi}
% =====================================================================

\subsection{SK Telecom: multi-vendor orchestration on KOREN}
\label{sec:sim}

Sim et al.~\cite{sim2023} describe the first documented commercial
deployment of a GS~QKD~018-compliant orchestrator, on the testbed of the
Korean research network KOREN (2020–2022). Two QKD operators run two
networks with devices and KMSs from different vendors, each with its own
controller (QKDNC); an orchestrator (QKDNO) coordinates them. Links use
decoy-state BB84 over spans between $39$ and $100$~km. The two networks meet
at \emph{border nodes} (e.g.\ Pangyo and Daejeon) that contain no QKD
modules but two interconnected KMSs, one per operator.

When two SAEs belong to different networks, the orchestrator computes the
end-to-end path with a \emph{path computation engine}, choosing \emph{in
advance} the border node to traverse; each controller then creates its own
service-link segment. To bind the segments the authors extend the model:
they introduce a \emph{multi-segment} application type (\texttt{MS}), with
an identifier \texttt{ms\_app\_id} shared across domains and
\texttt{server\_ms\_info} and \texttt{client\_ms\_info} describing the
external KMSs, plus an \texttt{app-registration-ms} RPC. Key transfer at
the border node uses the KMS–KMS interface of GS~QKD~020~\cite{qkd020}.
The case shows both the soundness of the 018 scheme and one of its limits:
the application types inherited from 015 (\texttt{CLIENT} and
\texttt{INTERNAL}) are not enough for multi-domain, and extensions are
needed. It is also a good illustration of the usefulness of the YANG
\texttt{identity} construct discussed in §\ref{sec:yang-anatomia}: a new
type can be added without rewriting the original model.

\subsection{Politecnico di Milano: hierarchical control with ONOS}

Vicario's thesis~\cite{vicario2025} builds an end-to-end prototype of
hierarchical control for metropolitan networks: the domain controllers are
based on ONOS and manage the nodes through the GS~QKD~015 model, while the
upward interface is GS~QKD~018, exposed via RESTCONF to a Python
orchestrator that aggregates the topology of several domains. It is a
useful implementation because it shows the whole 015~$\to$~018 stack in
open software, with the same YANG modules described in §\ref{sec:yang018}.

\subsection{MadQCI: the Madrid network}

MadQCI~\cite{madqci2024} is a real QKD network in Madrid with dozens of
modules from different vendors over more than $130$~km, integrated with the
production optical network. Control relies on an SD-QKD stack based on
GS~QKD~015, which is confirmed as the most widely adopted standard.
GS~QKD~018 is mentioned as a tool for inter-domain orchestration and
optical–quantum coordination, but the authors describe it as work still in
progress. This is an honest picture of the state of the art: domain control
is mature, orchestration much less so.
